
\medskip

\textbf{Les parties 1 et 2 sont indépendantes}

\bigskip

\textbf{Partie 1 : Distance de réaction}

\medskip

La distance de réaction d'un véhicule est la distance parcourue par ce véhicule entre l'instant où le conducteur voit un obstacle et l'instant où il appuie sur la pédale de frein.

On considère un conducteur en bonne santé.

La distance de réaction, en mètre, en fonction de la vitesse du véhicule est représentée par le graphique ci-dessous.


\begin{center}
\psset{xunit=0.1cm,yunit=0.2cm,arrowsize=2pt 3}
\begin{pspicture}(-10,-2.5)(120,40)
\multido{\n=0+2}{61}{\psline[linewidth=0.1pt](\n,0)(\n,35)}
\multido{\n=0+10}{13}{\psline[linewidth=0.5pt](\n,0)(\n,35)}
\multido{\n=0+1}{36}{\psline[linewidth=0.1pt](0,\n)(120,\n)}
\multido{\n=0+5}{8}{\psline[linewidth=0.5pt](0,\n)(120,\n)}
\uput[r](0,36.5){Distance (m)}\uput[u](105,0){Vitesse (km/h)}
\rput(60,38){Distance de réaction en fonction de la vitesse}
\psaxes[linewidth=1.25pt,Dx=10,Dy=5]{->}(0,0)(0,0)(120,35)
\psplot[plotpoints=500,linewidth=1.25pt,linecolor=blue]{0}{110}{0.27777 x mul}
\end{pspicture}
\end{center}

\medskip

\begin{enumerate}
\item Cette représentation graphique traduit-elle une situation de proportionnalité ?

Justifier la réponse.
\item Compléter, par lecture graphique, ce tableau.

\begin{center}
\begin{tabularx}{\linewidth}{|l| *{3}{>{\centering \arraybackslash}X|}}\hline
Vitesse (km/h)			&0 		&\ldots	&90\\ \hline
Distance de réaction (m)&\ldots	&15		&\ldots\\ \hline
\end{tabularx}
\end{center}

\end{enumerate}

\bigskip

\textbf{Partie 2 : Distance de freinage sur route sèche}

\medskip

La distance de freinage d'un véhicule est la distance parcourue par ce véhicule entre l'instant où le conducteur appuie sur la pédale de frein et l'instant où la voiture s'arrête complètement.

La distance de freinage en mètre, pour un véhicule en bon état, est déterminée en fonction de la vitesse du véhicule par la formule :

\begin{center}$d = \dfrac{v^2}{203,2}$ \:  où $v$ est la vitesse exprimée en km/h\end{center}

On utilise un tableur  pour calculer les distances de freinage en fonction de la vitesse :

\begin{center}
\begin{tabularx}{\linewidth}{|c|m{3.5cm}|*{3}{>{\centering \arraybackslash}X|}}\hline
&A&B&C&D\\ \hline
 1 &vitesse (km/h) &10 &20 &30\\ \hline
2 &distance de freinage (m)&&&\\ \hline
\end{tabularx}
\end{center}

\medskip

\begin{enumerate}
\item Recopier parmi les formules trois suivantes, celle qu'il faut saisir dans la cellule B2 puis étirer vers la droite:

\begin{center}\fbox{= 2*B1/203.2}\qquad  \fbox{= B1*B1/203.2} \qquad  \fbox{= B1+B1/203.2}\end{center}

\item Un véhicule roule à $90$ km/h. 

Montrer que sa distance de freinage est environ $40$~m.

\end{enumerate}

\bigskip

\textbf{Partie 3 : Distance d'arrêt sur route sèche}

\medskip

La distance d'arrêt d'un véhicule est la distance parcourue par ce véhicule entre l'instant où le conducteur voit un obstacle et l'instant où la voiture s'arrête complètement.

Distance d'arrêt = Distance de réaction $+$ Distance de freinage

Calculer la distance d'arrêt d'un véhicule roulant à $90$ km/h.

\bigskip

